\documentclass[11pt]{article}
\usepackage[margin=.7in,top=.55in,bottom=.55in]{geometry}
\usepackage{amsmath,amssymb,amsthm,hyperref}
\renewcommand{\normalsize}{\fontsize{10}{12}\selectfont}
\DeclareMathSizes{10}{10}{8}{6}
\hypersetup{colorlinks=true,linkcolor=blue,urlcolor=blue,citecolor=blue}
\setlength{\emergencystretch}{2em}
\newtheorem{theorem}{Theorem}
\newtheorem{lemma}{Lemma}
\newcommand{\A}{\mathcal A}
\newcommand{\D}{\overline{\mathbb D^*}}
\begin{document}
\begin{center}
{\Large\bfseries A continuous critical-branch obstruction\\to C*-algebraic Sendov conjectures}\\[.5em]
Research proof; formalization draft\\
October 26, 2026
\end{center}
\begin{abstract}
We give a cubic in the commutative unital C*-algebra
$C([0,1],\mathbb C)$ whose roots have norm at most $9/10$, whose derivative
splits completely into two distinct continuous critical branches of norm
less than one, and whose critical branches are positive continuous convex
combinations of the original roots. Nevertheless, every critical element
has distance greater than $21/20$ from one fixed root. This contradicts
Conjecture~3 of Krishna's published C*-algebraic Sendov paper, originally
Conjecture~2.5 of its preprint, and its completely split derivative variant
(published Conjecture~4). The obstruction concerns the existence of one
global continuous critical element. It does not contradict scalar Sendov.
The two source formulations are treated as one combined outcome.
\end{abstract}

\section{Source statement and the counterexample}
For a unital C*-algebra, Krishna \cite[Definition~2]{K26} defines
\[
 \D(a,r)=\{z\in\A:(z-a)(z-a)^*\le\sqrt r\,1\}.
\]
Only $r=1$ occurs here. In $\A=C([0,1],\mathbb C)$, equipped with its
pointwise operations and supremum norm, this condition is precisely
$|z(t)-a(t)|\le1$ for every $t\in[0,1]$.
For $p(z)=\prod_{j=0}^2(z-a_j)$ the source defines $p'(z)$ by the sum of
the three products obtained by deleting one factor, preserving order.
The algebra used here is commutative, so this is also ordinary formal
polynomial differentiation.

Published Conjecture~3 \cite[p.~60]{K26}, corresponding to preprint
Conjecture~2.5 \cite{K22} and MathDB~353909, assumes that all $a_j$
belong to $\D(0,1)$, that $p'$ has two roots $b_k\in\D(0,1)$, and
that each $b_k$ has a representation
\begin{equation}\label{convex}
 b_k=\sum_{j=0}^2\omega_{kj}a_j,\qquad
 \sum_{j=0}^2\omega_{kj}=1,\qquad \omega_{kj}\ge0\text{ in }\A.
\end{equation}
It asserts that for each $a_j$ there is a zero $b$ of $p'$ in
$\D(a_j,1)$. Published Conjecture~4 replaces the critical-root and convexity
hypotheses with complete linear factorization of the derivative.

\begin{theorem}\label{main}
There exist three elements $a_0,a_1,a_2\in C([0,1],\mathbb C)$ of norm
at most $9/10$ and two distinct elements $B_+,B_-$ such that:
\begin{enumerate}
\item $p$ is monic of degree three and
$p'(z)=3(z-B_+)(z-B_-)$ for every $z\in\A$;
\item $\|B_\pm\|<9/10$, and both $B_\pm$ satisfy \eqref{convex} with
strictly positive continuous real-valued coefficients;
\item every $b\in\A$ satisfying $p'(b)=0$ is either $B_+$ or $B_-$, and
$\|a_0-b\|>21/20$.
\end{enumerate}
Consequently the conclusions of both source formulations fail.
\end{theorem}

Put $u(t)=e^{\pi it}$ and define
\begin{equation}\label{roots}
 a_0=\frac9{10},\qquad a_1(t)=\frac{27}{40}u(t),\qquad
 a_2(t)=-\frac{27}{40}u(t).
\end{equation}
The three roots are pairwise distinct at every point: $|a_0|=9/10$,
$|a_1|=|a_2|=27/40$, and $u$ is nowhere zero.
In particular all roots lie strictly inside the required unit disc.

\section{Continuous critical branches and complete splitting}
Use the principal complex square root to set
\begin{equation}\label{branches}
 h(t)=1+\frac{16}{27}\overline{u(t)}^{\,2},\qquad
 q(t)=\frac{3\sqrt3}{4}u(t)\sqrt{h(t)},\qquad
 B_\pm(t)=\frac3{10}(1\pm q(t)).
\end{equation}
Since $\operatorname{Re}h\ge11/27>0$, these functions are continuous.
Also $u\overline u=1$, so direct multiplication gives
\begin{equation}\label{qsquare}
 q^2=1+\frac{27}{16}u^2.
\end{equation}
The right-hand side has modulus at least $11/16$. Thus $q$ is nowhere
zero, and the two branches are distinct everywhere. Furthermore,
$|q|^2=|q^2|\le43/16<4$. Consequently
\[
 |B_\pm(t)|\le\frac3{10}(1+|q(t)|)<\frac9{10}.
\]
Compactness makes this a strict supremum-norm bound as well.

Expanding the source's literal deleted-factor derivative yields
\begin{align*}
 p'(z)&=(z-a_1)(z-a_2)+(z-a_0)(z-a_2)+(z-a_0)(z-a_1)\\
 &=3z^2-\frac95z-\frac{729}{1600}u^2
 =3(z-B_+)(z-B_-).
\end{align*}
The last equality follows from \eqref{qsquare}. It holds both as an identity
in $\A[z]$ and after evaluation at every $z\in\A$.

\begin{lemma}\label{allcritical}
The only zeros of $p'$ in $\A$ are $B_+$ and $B_-$.
\end{lemma}
\begin{proof}
If $p'(b)=0$, then at each $t$ the scalar factorization gives
$b(t)=B_+(t)$ or $b(t)=B_-(t)$. The continuous function
\[
 F(t)=\frac{(10/3)b(t)-1}{q(t)}
\]
therefore takes values in $\{-1,1\}$. The interval is connected, so $F$
is constant. Its two possible values give the two asserted global branches.
This argument classifies all continuous critical elements despite the zero
divisors in the function algebra.
\end{proof}

\section{Positive convex weights and the endpoint obstruction}
Fix either branch $B$. As the three scalar roots are simple at every $t$,
no critical point can equal any $a_j(t)$: evaluating $p'_t$ at $a_j(t)$
gives the product of the other two nonzero differences. Thus
\[
 r_j(t)=\frac1{|B(t)-a_j(t)|^2},\qquad
 R(t)=\sum_{j=0}^2r_j(t),\qquad
 \omega_j(t)=\frac{r_j(t)}{R(t)}
\]
are continuous, $r_j,R,\omega_j$ are strictly positive, and $\sum_j\omega_j=1$.
Pointwise division of $p'(B)=0$ by the nowhere-zero product $p(B)$ gives
$\sum_j(B-a_j)^{-1}=0$. Conjugating this equality gives
$\sum_j r_j(B-a_j)=0$, hence $B=\sum_j\omega_j a_j$.
Real-valued nonnegative continuous functions are positive elements of $\A$,
so these are exactly the coefficients required in \eqref{convex}.

At the endpoints put $Q=\sqrt{43}/4$. Formula \eqref{branches} gives
\[
 q(0)=Q>\frac32,\qquad q(1)=-Q,\qquad
 B_-(0)=B_+(1)=\frac3{10}(1-Q).
\]
Therefore
\begin{equation}\label{gap}
 |a_0(0)-B_-(0)|=|a_0(1)-B_+(1)|
 =\frac35+\frac{3\sqrt{43}}{40}>\frac{21}{20}>1.
\end{equation}
Each branch fails the unit-disc condition at one endpoint. Taking the
supremum and applying Lemma~\ref{allcritical} proves Theorem~\ref{main}.
The actual endpoint distance is approximately $1.0918078893$.

\section{Scope and verification status}
The construction retains complete derivative splitting, strictly interior
roots and critical roots, and positive continuous convex weights. The
failure is the simultaneous requirement of continuity and one uniform
nearby critical element for a prescribed root. The underlying connectedness
and square-root mechanism is classical and receives no separate novelty
claim. The scalar pointwise Sendov statement is compatible with the example.

The commutative source formulation, MathDB~353908, also assumes density of
invertibles. That additional hypothesis is not part of the formal target
here; the reported source refutation is MathDB~353909 and the published
completely split derivative variant, counted together. Formal implementation
and independent bounded prior-art review are in progress. No acceptance,
complete formal certificate, or absolute priority is asserted in this draft.

\begin{thebibliography}{9}
\bibitem{K26} K.~M.~Krishna,
\emph{C*-algebraic Gauss--Lucas theorem and C*-algebraic Sendov's conjecture},
Vestnik KRAUNC. Fiz.-Mat. Nauki \textbf{54}(1) (2026), 56--63.
\href{https://doi.org/10.26117/2079-6641-2026-54-1-56-63}{doi:10.26117/2079-6641-2026-54-1-56-63}.
\bibitem{K22} K.~M.~Krishna, same title,
\href{https://arxiv.org/abs/2203.06916v1}{arXiv:2203.06916v1} (2022).
\end{thebibliography}
\end{document}
